% !PRINTED 616
\weberpage{616}
%<<CONT
\weq{8}{\pi \equiv \beta \quad (\mathrm{mod}\ \mathfrak{m})}
\gesperrt{genügen, für die}
\weq{9}{\mathfrak{p} = \mathfrak{a}\pi}
\gesperrt{ein Primideal wird.}

Nimmt man $\mathfrak{a} = 1$, so wird $\pi$ eine ganze Zahl, also eine in $\Omega$ existierende Primzahl. Der Satz lautet für diesen besonderen Fall:

3. \gesperrt{Es gibt im Körper $\Omega$ unendlich viele existierende Primzahlen $\pi$, deren Norm eine natürliche Primzahl ist, die nach einem beliebigen Modul $\mathfrak{m}$ mit einer beliebigen zu $\mathfrak{m}$ teilerfremden Zahl $\beta$ kongruent sind.}

Nehmen wir für den Modul $\mathfrak{m}$ eine ganze Zahl $\mu$ in $\Omega$, so erhalten wir folgenden Satz:

4. \gesperrt{Ist $\mu$ eine ganze Zahl in $\Omega$, so sind in der Linearform}
\[
\mu\xi + \beta,
\]
\gesperrt{in der $\xi$ die ganzen Zahlen von $\Omega$ durchläuft und $\beta$ relativ prim zu $\mu$ ist, unendlich viele Primzahlen enthalten.}

Man hat hier eine schöne Verallgemeinerung des Satzes von den Primzahlen in arithmetischen Progressionen (Bd.\ II, \S\,198)\footnote{Für den Fall $\Delta = -4$ ist dieser Beweis gegeben von \gesperrt{H. Weber}, Crelles Journ., Bd.\ 129 (Dirichlet-Band), für $\Delta = -3$ in einer Straßburger Dissertation von \gesperrt{H. Bresslau} (Straßburg 1907).}.

\begin{center}
\textbf{\S\,168. Reduktion der Klassengleichung in den Kreisteilungskörpern.}
\end{center}

Wir definieren eine weitere Gruppe $A$ durch folgende Bestimmung:

Es seien $Q$, $m$ zwei ganze rationale Zahlen, deren Primfaktoren wir in den Exkludenten $S$ aufnehmen. $O$ sei wie in \S\,98 die Gruppe der Zahlen in $\Omega$, und $O'$ die Gruppe der Zahlen der Ordnung $[Q]$, und $A$ bestehe aus allen ganzen und gebrochenen Zahlen $\alpha$ in $O'$, die der Kongruenzbedingung
\weq{1}{N(\alpha) \equiv 1 \quad (\mathrm{mod}\ m)}
genügen. Zunächst haben wir die Klassenzahl
\weq{2}{h = (\overline{O}, \overline{E}A) = (\overline{O}, \overline{E}O')\,(\overline{E}O', \overline{E}A)}
%CONT>>
